\section{Introduction}
\label{sec:introduction}

For a smooth cubic surface whose twenty-seven lines have the full
incidence-compatible Galois group \(\WE\), the algebraic Brauer quotient may
be trivial over the base yet acquire 2-torsion after a finite extension.  How
small can that extension be?  We answer this for the exact surface \(Y/\Q\)
in \cref{eq:cubic}.  The companion theorem package
\code{henon\_mu3\_yukawa\_line\_field}, included in the anonymous supplement,
proves that the common normal field \(K\) of its lines satisfies
\[
  \Gal(K/\Q)\cong\WE,\qquad |\WE|=51840.
\]
The degree-\(27\) field \(E\) of one line is non-Galois and differs from
\(K\).  These are frozen inputs: we locate the first jump inside \(K\), not
recompute the line scheme.

Nonzero 2-primary cohomology has two classical branches.  A quotient
\(\Z/2\) places the Galois image in a conjugate of the double-six stabilizer
\(U_1\cong S_6\times C_2\), of index \(36\); a quotient
\((\Z/2)^2\) places it in a conjugate of \(U_3\), of index \(720\).
Finding one invariant double-six does not eliminate the second branch.  The
complete
Swinnerton-Dyer--Elsenhans--Jahnel classification supplies this universal
dichotomy \cite{SwinnertonDyer1993,ElsenhansJahnel2010Brauer}.

The result combines that classification with exact configuration arithmetic
and a written descent argument.  For an abelian group \(A\), write \(A[2]\)
for its subgroup annihilated by \(2\).  For a finite extension \(L/\Q\), set
\begin{equation}
 G_L=\Gal(\Qbar/L),\qquad
 H_L=\im(G_L\to\WE),\qquad
 N_L=\ker(G_L\to H_L).
\label{eq:intro-GL-HL-NL}
\end{equation}
The cohomological passage from \(G_L\) to \(H_L\) is proved in
\cref{lem:inflation,eq:brauer-picard-bridge}, applied with \(k=L\).

\begin{theorem}[Main theorem]
\label{thm:main}
Let \(Y/\Q\) be the cubic surface in \cref{eq:cubic}, and let \(K/\Q\) be
the common normal field of its twenty-seven lines.  For every finite extension
\(L/\Q\),
\begin{equation}
 \bigl(\Brq{L}\bigr)[2]\ne0
 \quad\Longrightarrow\quad
 36\mid[K\cap L:\Q]\mid[L:\Q].
\label{eq:intro-divisibility}
\end{equation}
Consequently the least possible degree is \(36\).  If
\([L:\Q]=36\) and the left side of \cref{eq:intro-divisibility} is nonzero,
then, inside the fixed algebraic closure,
\[
 L=K\cap L=K^{wU_1w^{-1}}
\]
for some \(w\in\WE\), and this subgroup stabilizes a unique embedded
double-six.  The thirty-six embedded
fields so obtained form one \(\Q\)-isomorphism type.

For a selected double-six \(D=\{\cE,\cG\}\), put
\(\FD=K^{U_1}\).  There are exact invariants
\(\theta_D,\delta_D\in K\) such that
\[
 \Q(\theta_D)=\FD=\Q(\delta_D),\qquad
 \FDp=K^{U_1^+}=\FD(\sqrt{\delta_D}),
\]
where \(U_1^+\cong S_6\) preserves the two sixers separately.  Moreover,
\[
 \Brq{\Q}=0,\qquad \Brq{\FD}\cong\Z/2.
\]
There is a unique normalized determinant-defined quartic \(Q_D\) satisfying
\[
 \divv_{Y_{\FD}}(Q_D)=\cE+\cG,
\]
and the quaternion
\[
 \mathcal A_D=(\delta_D,Q_D/u_0^4)
\]
is unramified and represents the unique nonzero element of \(\Brq{\FD}\).
\end{theorem}

The contributions are:
\begin{enumerate}[leftmargin=*]
\item the full divisibility chain and equality classification through both
      \(U_1/U_3\) branches;
\item exact degree-\(36\) fields, their common normal closure \(K\), and the
      oriented quadratic extension;
\item a determinant quartic whose degree-exhausted divisor and integral
      Picard cocycle prove that \((\delta_D,Q_D/u_0^4)\) is unramified and
      nonzero.
\end{enumerate}

\begin{figure}[t]
\centering
\setlength{\fboxsep}{8pt}
\fbox{\begin{minipage}{0.92\textwidth}
\centering
\small
\begin{tabularx}{\linewidth}{@{}>{\centering\arraybackslash}p{0.16\linewidth}
 >{\centering\arraybackslash}X>{\centering\arraybackslash}X@{}}
\toprule
branch & forced Galois image & degree consequence\\
\midrule
\(\Z/2\)
  & \(H_L\subseteq wU_1w^{-1}\), \([\WE:U_1]=36\)
  & equality occurs at \(\FD=K^{U_1}\)\\[4pt]
\((\Z/2)^2\)
  & \(H_L\subseteq wU_3w^{-1}\), \([\WE:U_3]=720\)
  & excluded when \([L:\Q]=36\)\\
\bottomrule
\end{tabularx}

\vspace{5pt}
\[
 K\supset \FDp=K^{U_1^+}=\FD(\sqrt{\delta_D})
 \supset \FD=K^{U_1}\supset\Q.
\]
\begin{align*}
 \FD&\quad\Longrightarrow\quad
 (\delta_D,Q_D/u_0^4)\in\Br(Y_{\FD}),\\
 [K\cap L:\Q]&=[\WE:H_L],
 \qquad 36\mid[K\cap L:\Q]\mid[L:\Q].
\end{align*}
\end{minipage}}
\caption{Both classification branches force divisibility by \(36\), but only
\(U_1\) attains equality.  Its fixed field, orientation extension, and
quaternion form the attaining package; no local evaluation is asserted.}
\label{fig:field-subgroup}
\end{figure}

Invariant double-sixes, orientation fields, and cyclic algebras are
established mechanisms
\cite{ElsenhansJahnel2010Brauer,ElsenhansJahnel2010DoubleSix}; degree-\(36\)
resolvers and their moduli context are also prior art
\cite{ElsenhansJahnel2012OrderThree,FarbWolfson2019}.  Our claim is
instance-specific: the bounded 2026-08-15 screen did not locate a prior exact
computation of the degree-36 double-six field, its orientation square, and
the determinant-defined quaternion generator for this frozen Yukawa surface.

The proof is organized around the path visible in
\cref{fig:field-subgroup}.  \Cref{sec:setup-sources} fixes the surface,
notation, sources, and evidence contract.  \Cref{sec:minimal-degree} proves
the universal minimum.  \Cref{sec:double-six-fields,sec:brauer-jump}
identify the fields and cohomological jump.  The determinant quartic is
constructed in \cref{sec:canonical-quartic}; its divisor and quaternion class
are proved in \cref{sec:quaternion}.  Exact replay and limitations appear in
\cref{sec:replay-scope}.  We compute no local evaluation and draw no
conclusion about rational points, the Hasse principle, weak approximation, or
a Brauer--Manin obstruction.
